% ============================================================
%  CAHIER DES CHARGES — OUTIL D'AUDIT DE CONFIGURATION
%  Cisco IOS / IOS-XE Security Assessment (OFFLINE)
% ============================================================
\documentclass[12pt, a4paper]{report}

% ---------- Packages ----------
\usepackage[T1]{fontenc}
\usepackage[utf8]{inputenc}
\usepackage[french]{babel}
\usepackage{lmodern}
\usepackage{geometry}
\usepackage{graphicx}
\usepackage{xcolor}
\usepackage{amsmath}
\usepackage{amssymb}
\usepackage{array}
\usepackage{multirow}
\usepackage{booktabs}
\usepackage{longtable}
\usepackage{tabularx}
\usepackage{colortbl}
\usepackage{float}
\usepackage{caption}
\usepackage{setspace}
\usepackage{parskip}
\usepackage{microtype}
\usepackage{enumitem}
\usepackage{fancyhdr}
\usepackage{titlesec}
\usepackage{titletoc}
\usepackage{mdframed}
\usepackage{listings}
\usepackage{tikz}
\usetikzlibrary{shapes.geometric, arrows, positioning, shadows, calc}
% --- hyperref MUST come after most packages ---
\usepackage{hyperref}
% --- tcolorbox and fontawesome5 AFTER hyperref ---
\usepackage{tcolorbox}
\usepackage{fontawesome5}

% ---------- Geometry ----------
\geometry{
  top=2.5cm, bottom=2.5cm,
  left=2.8cm, right=2.8cm,
  headheight=15pt
}

% ---------- Colors ----------
\definecolor{ciscoBlue}{HTML}{049FD9}
\definecolor{ciscoDark}{HTML}{005073}
\definecolor{ciscoGray}{HTML}{58585B}
\definecolor{pastelBlue}{HTML}{EBF5FB}
\definecolor{pastelGray}{HTML}{F5F5F5}
\definecolor{accentRed}{HTML}{E74C3C}
\definecolor{accentGreen}{HTML}{27AE60}
\definecolor{accentOrange}{HTML}{F39C12}
\definecolor{tableHeader}{HTML}{005073}
\definecolor{tableRowA}{HTML}{EBF5FB}
\definecolor{tableRowB}{HTML}{FFFFFF}
\definecolor{rulePass}{HTML}{D5F5E3}
\definecolor{ruleFail}{HTML}{FADBD8}
\definecolor{ruleWarn}{HTML}{FDEBD0}

% ---------- Hyperref ----------
\hypersetup{
  colorlinks=true,
  linkcolor=ciscoDark,
  urlcolor=ciscoBlue,
  citecolor=ciscoDark,
  pdftitle={Cahier des Charges - Audit Cisco IOS-XE},
  pdfauthor={Equipe Projet},
  pdfsubject={Security Configuration Assessment - Cisco IOS/IOS-XE}
}

% ---------- Section Styles ----------
\titleformat{\chapter}[block]
  {\color{ciscoDark}\huge\bfseries}
  {\thechapter}{1em}{}
  [\vspace{2pt}\color{ciscoBlue}\rule{\textwidth}{2pt}]

\titleformat{\section}
  {\color{ciscoDark}\Large\bfseries}
  {\thesection}{1em}{}
  [\vspace{1pt}\color{ciscoBlue}\rule{\textwidth}{0.8pt}]

\titleformat{\subsection}
  {\color{ciscoDark}\large\bfseries}
  {\thesubsection}{1em}{}

\titleformat{\subsubsection}
  {\color{ciscoGray}\normalsize\bfseries}
  {\thesubsubsection}{1em}{}

% ---------- Header / Footer ----------
\pagestyle{fancy}
\fancyhf{}
\fancyhead[L]{\small\color{ciscoGray}\textit{Cahier des Charges — Audit Cisco IOS/IOS-XE}}
\fancyhead[R]{\small\color{ciscoGray}\textit{Version 1.0 — \today}}
\fancyfoot[C]{\color{ciscoGray}\small\thepage}
\renewcommand{\headrulewidth}{0.4pt}
\renewcommand{\footrulewidth}{0.4pt}
\renewcommand{\headrule}{\hbox to\headwidth{\color{ciscoBlue}\leaders\hrule height \headrulewidth\hfill}}

% ---------- Listings (code) ----------
\lstdefinestyle{iosStyle}{
  backgroundcolor=\color{pastelGray},
  basicstyle=\ttfamily\small,
  breaklines=true,
  captionpos=b,
  commentstyle=\color{ciscoGray},
  keywordstyle=\color{ciscoDark}\bfseries,
  stringstyle=\color{ciscoBlue},
  frame=single,
  rulecolor=\color{ciscoBlue},
  numbers=left,
  numberstyle=\tiny\color{ciscoGray},
  numbersep=5pt,
  showstringspaces=false,
  tabsize=2
}
\lstset{style=iosStyle}

% ---------- Custom Boxes ----------
\tcbuselibrary{skins, breakable, listings}

\newtcolorbox{infobox}[1][]{
  enhanced,
  colback=pastelBlue,
  colframe=ciscoBlue,
  fonttitle=\bfseries\color{white},
  coltitle=white,
  attach boxed title to top left={yshift=-2mm, xshift=4mm},
  boxed title style={colback=ciscoBlue},
  title=Information,
  #1
}

\newtcolorbox{rulebox}[2][]{
  enhanced,
  colback=pastelGray,
  colframe=ciscoDark,
  fonttitle=\bfseries\color{white},
  coltitle=white,
  attach boxed title to top left={yshift=-2mm, xshift=4mm},
  boxed title style={colback=ciscoDark},
  title=#2,
  #1
}

\newtcolorbox{warnbox}[1][]{
  enhanced,
  colback=ruleWarn,
  colframe=accentOrange,
  fonttitle=\bfseries,
  title=\faExclamationTriangle\ Attention,
  #1
}

% ============================================================
\begin{document}

% ============================================================
%  PAGE DE TITRE
% ============================================================
\begin{titlepage}
  \thispagestyle{empty}

  % ---- Top header bar ----
  \noindent
  \begin{tikzpicture}[remember picture, overlay]
    \fill[ciscoDark] (current page.north west)
      rectangle ([yshift=-2.2cm]current page.north east);
    \fill[ciscoBlue] ([yshift=-2.2cm]current page.north west)
      rectangle ([yshift=-2.45cm]current page.north east);
  \end{tikzpicture}

  % ---- Logo (top-left inside header) ----
  \begin{flushright}
    \vspace*{0.3cm}
    \includegraphics[height=3cm]{logo.png}
  \end{flushright}

  \vspace{1.6cm}

  % ---- Main title block ----
  \begin{center}
    {\fontsize{9pt}{11pt}\selectfont\color{ciscoGray}\MakeUppercase{%
      Document de Spécification Technique — Confidentiel}}
    \\[0.4cm]
    \color{ciscoDark}\rule{0.9\textwidth}{0.6pt}\\[0.6cm]
    {\fontsize{28pt}{32pt}\selectfont\bfseries\color{ciscoDark}
      CAHIER DES CHARGES}\\[0.35cm]
    {\fontsize{16pt}{20pt}\selectfont\bfseries\color{ciscoBlue}
      Outil d'Audit de Configuration Sécurité}\\[0.25cm]
    {\fontsize{13pt}{16pt}\selectfont\bfseries\color{ciscoGray}
      Cisco IOS\,/\,IOS-XE — Évaluation OFFLINE}\\[0.5cm]
    \color{ciscoDark}\rule{0.9\textwidth}{0.6pt}
  \end{center}

  \vspace{0.9cm}

  % ---- Synopsis box ----
  \begin{center}
  \begin{tcolorbox}[
    enhanced,
    colback=pastelBlue,
    colframe=ciscoDark,
    leftrule=5pt, rightrule=0.4pt, toprule=0.4pt, bottomrule=0.4pt,
    width=0.88\textwidth,
    arc=2mm,
    left=8pt, right=8pt, top=6pt, bottom=6pt
  ]
    \small\color{ciscoDark}
    \textbf{Objet :} Ce document définit les exigences fonctionnelles et techniques de
    l'outil d'audit de configuration sécurité pour équipements Cisco IOS\,/\,IOS-XE.
    L'analyse est réalisée \textbf{hors-ligne} à partir du fichier
    \texttt{show running\-config}, en conformité avec le guide officiel Cisco
    \textit{«Harden IOS/IOS-XE Devices»}.
    Les résultats sont restitués sous forme de rapport \textbf{JSON} et \textbf{HTML}.
  \end{tcolorbox}
  \end{center}

  \vspace{1.1cm}

  % ---- Document info table ----
  \begin{center}
  {%
    \arrayrulecolor{ciscoBlue}%
    \setlength{\arrayrulewidth}{1.2pt}%
    \begin{tabular}{|>{\bfseries\color{ciscoDark}\small}p{3.8cm}
                    |>{\small\color{black}}p{8.5cm}|}
      \hline
      \rowcolor{tableHeader}
      \color{white}\small\bfseries Champ & \color{white}\small\bfseries Valeur \\
      \hline
      \rowcolor{tableRowA}
      Projet      & Automatisation d'Audit de Configuration \\
      \rowcolor{tableRowB}
      Module      & Cisco IOS\,/\,IOS-XE Security Hardening \\
      \rowcolor{tableRowA}
      Version     & 1.0 \\
      \rowcolor{tableRowB}
      Date        & \today \\
      \rowcolor{tableRowA}
      Référentiel & Cisco Guide — Harden IOS/IOS-XE Devices \\
      \rowcolor{tableRowB}
      Statut      & \textcolor{accentGreen}{\bfseries En cours de rédaction} \\
      \hline
    \end{tabular}%
    \arrayrulecolor{black}%
  }
  \end{center}

  \vfill

  % ---- Bottom bar ----
  \noindent
  \begin{tikzpicture}[remember picture, overlay]
    \fill[ciscoBlue] (current page.south west)
      rectangle ([yshift=0.35cm]current page.south east);
    \fill[ciscoDark] ([yshift=0.35cm]current page.south west)
      rectangle ([yshift=1.1cm]current page.south east);
    \node[anchor=center, text=white, font=\small\itshape]
      at ([yshift=0.72cm]current page.south)
      {Confidentiel — Usage Interne Uniquement};
  \end{tikzpicture}
  \vspace{0.8cm}
\end{titlepage}

% ============================================================
%  TABLE DES MATIÈRES
% ============================================================
\tableofcontents
\clearpage

% ============================================================
\chapter{Contexte et Présentation du Projet}
% ============================================================

\section{Contexte Général}

Dans un contexte où la sécurité des infrastructures réseau est devenue un enjeu stratégique majeur, les organisations doivent s'assurer que leurs équipements Cisco IOS/IOS-XE sont configurés conformément aux meilleures pratiques de durcissement (hardening). Ces vérifications, réalisées manuellement par des auditeurs, sont chronophages, sujettes à l'erreur humaine et difficiles à répliquer à grande échelle.

Ce document constitue le \textbf{cahier des charges} de l'outil d'audit de configuration sécurité OFFLINE pour équipements Cisco IOS et IOS-XE. Il décrit de manière exhaustive les règles de conformité à implémenter, la logique d'évaluation attendue, les livrables techniques et les contraintes de développement.

\begin{infobox}[title=Périmètre de ce document]
  Ce cahier des charges se concentre exclusivement sur le module \textbf{Cisco IOS/IOS-XE}. Les modules Windows et Linux (Ubuntu) feront l'objet de documents distincts. La référence technique principale est le guide Cisco \textit{«Harden IOS / IOS-XE Devices»} disponible à l'adresse : \url{https://www.cisco.com/c/en/us/support/docs/ip/access-lists/13608-21.html}
\end{infobox}

\section{Problématique}

Les équipements réseau Cisco constituent le cœur des infrastructures IT. Une mauvaise configuration peut exposer le réseau à de nombreuses vulnérabilités :
\begin{itemize}[leftmargin=1.5cm]
  \item Accès non autorisé aux interfaces de management (SSH, Telnet, SNMP)
  \item Absence de mécanismes d'authentification robustes (AAA, TACACS+)
  \item Exposition aux attaques de déni de service (DoS) par absence de CoPP ou de rate-limiting
  \item Usurpation d'adresse IP par absence de Unicast RPF ou d'IP Source Guard
  \item Absence de traçabilité (logs, NTP, NetFlow)
  \item Protocoles de routage non sécurisés (absence d'authentification MD5)
\end{itemize}

\section{Objectifs du Projet}

\begin{enumerate}[leftmargin=1.5cm, label=\textbf{OBJ-\arabic*.}]
  \item \textbf{Importation OFFLINE} : Analyser les fichiers \texttt{show running-config} exportés sans connexion directe aux équipements.
  \item \textbf{Évaluation de règles} : Vérifier chaque règle de conformité et retourner un statut \texttt{PASS / FAIL / NOT FOUND / NOT APPLICABLE}.
  \item \textbf{Génération de rapport} : Produire un rapport lisible au format JSON et HTML avec score global et détail par règle.
  \item \textbf{Extensibilité} : Permettre l'ajout de nouvelles règles via des fichiers YAML sans modifier le code Python.
  \item \textbf{Couverture} : Couvrir les trois plans de l'architecture Cisco : Management Plane, Control Plane, Data Plane.
\end{enumerate}

\section{Périmètre Technique}

\begin{table}[H]
  \centering
  \caption{Périmètre technique du module Cisco}
  \begin{tabularx}{\textwidth}{>{\bfseries\color{ciscoDark}}l X}
    \toprule
    \rowcolor{tableHeader}\color{white}\textbf{Élément} & \color{white}\textbf{Description} \\
    \midrule
    \rowcolor{tableRowA} Équipements visés & Routeurs et switches Cisco IOS / IOS-XE \\
    \rowcolor{tableRowB} Source de données & Fichier texte brut : \texttt{show running-config} \\
    \rowcolor{tableRowA} Mode de fonctionnement & OFFLINE (aucune connexion réseau requise) \\
    \rowcolor{tableRowB} Référentiel & Cisco Harden IOS/IOS-XE Devices Guide \\
    \rowcolor{tableRowA} Langage & Python 3.10+ \\
    \rowcolor{tableRowB} Format des règles & YAML \\
    \rowcolor{tableRowA} Sorties & \texttt{report.json} + \texttt{report.html} \\
    \bottomrule
  \end{tabularx}
\end{table}

% ============================================================
\chapter{Architecture Logicielle}
% ============================================================

\section{Vue d'Ensemble}

L'outil est structuré autour de quatre composants principaux :

\begin{enumerate}[leftmargin=1.5cm]
  \item \textbf{Parseur Cisco} (\texttt{cisco\_parser.py}) : Lit et structure le contenu du \texttt{show running-config}.
  \item \textbf{Moteur d'évaluation} (\texttt{evaluator.py}) : Applique les règles YAML sur la configuration parsée.
  \item \textbf{Base de règles} (\texttt{rules/cisco\_rules.yaml}) : Définit toutes les règles de conformité.
  \item \textbf{Générateur de rapports} (\texttt{reporter.py}) : Produit les fichiers \texttt{report.json} et \texttt{report.html}.
\end{enumerate}

\section{Arborescence du Projet}

\begin{lstlisting}[language=bash, caption={Structure des fichiers du module Cisco}]
config-assessment-tool/
|-- cisco/
|   |-- cisco_parser.py          # Parseur du running-config
|   |-- evaluator.py             # Moteur de regles
|   |-- reporter.py              # Generateur de rapports
|   |-- rules/
|   |   |-- management_plane.yaml
|   |   |-- control_plane.yaml
|   |   |-- data_plane.yaml
|   |-- samples/
|   |   |-- sample_running_config.txt
|   |-- tests/
|       |-- test_cisco_parser.py
|       |-- test_evaluator.py
|-- output/
|   |-- report.json
|   |-- report.html
|-- README.md
\end{lstlisting}

\section{Statuts d'Évaluation}

Chaque règle peut retourner l'un des quatre statuts suivants :

\begin{table}[H]
  \centering
  \caption{Statuts de conformité}
  \begin{tabularx}{\textwidth}{>{\bfseries}l l X}
    \toprule
    \rowcolor{tableHeader}\color{white}\textbf{Statut} & \color{white}\textbf{Couleur} & \color{white}\textbf{Signification} \\
    \midrule
    \rowcolor{rulePass} PASS & Vert & La règle est respectée dans la configuration. \\
    \rowcolor{ruleFail} FAIL & Rouge & La règle n'est pas respectée ou la commande est absente. \\
    \rowcolor{ruleWarn} NOT FOUND & Orange & L'élément à vérifier est absent du fichier de configuration. \\
    \rowcolor{tableRowB} NOT APPLICABLE & Gris & La règle ne s'applique pas au contexte de cet équipement. \\
    \bottomrule
  \end{tabularx}
\end{table}

\section{Format d'une Règle YAML}

Chaque règle est définie dans un fichier YAML selon le format suivant :

\begin{lstlisting}[language=bash, caption={Exemple de règle YAML pour SSHv2}]
- id: CISCO-MGT-SSH-001
  section: Management Plane
  subsection: Secure Interactive Management Sessions
  title: "SSHv2 doit etre active"
  description: >
    SSHv2 est la version securisee du protocole SSH. SSHv1 presente
    des vulnerabilites connues. La commande 'ip ssh version 2' doit
    etre presente dans le running-config.
  check_type: regex_present
  pattern: "ip ssh version 2"
  remediation: >
    Executer : ip ssh version 2
  severity: HIGH
  reference: "Cisco Harden IOS-XE - SSHv2"
\end{lstlisting}

% ============================================================
\chapter{Management Plane — Règles de Conformité}
% ============================================================

Le \textbf{Management Plane} regroupe toutes les fonctions permettant de gérer, surveiller et administrer l'équipement réseau. C'est le plan le plus exposé car il concentre les accès d'administration.

% ---- 3.1 ----
\section{General Management Plane Hardening}

\subsection{Password Management}

\subsubsection{Enhanced Password Security}

\begin{rulebox}{CISCO-MGT-PWD-001 — Enhanced Password Security}
  \textbf{Objectif :} S'assurer que les mots de passe sont stockés sous forme chiffrée (type 5 ou type 8/9) et non en clair ou en type 7 (réversible).

  \textbf{Commandes à vérifier :}
  \begin{lstlisting}
service password-encryption
enable secret <hash>   ! Type 5 minimum recommande
  \end{lstlisting}

  \textbf{Logique de vérification :}
  \begin{itemize}
    \item \texttt{PASS} : \texttt{service password-encryption} présent ET \texttt{enable secret} (non \texttt{enable password}) utilisé.
    \item \texttt{FAIL} : \texttt{enable password} en clair détecté OU absence de \texttt{service password-encryption}.
  \end{itemize}

  \textbf{Sévérité :} \textcolor{accentRed}{\textbf{CRITIQUE}}
\end{rulebox}

\subsubsection{Login Password Retry Lockout}

\begin{rulebox}{CISCO-MGT-PWD-002 — Login Password Retry Lockout}
  \textbf{Objectif :} Bloquer un compte après N tentatives d'authentification échouées pour prévenir les attaques par force brute.

  \textbf{Commandes à vérifier :}
  \begin{lstlisting}
aaa local authentication attempts max-fail <N>
  \end{lstlisting}

  \textbf{Logique de vérification :}
  \begin{itemize}
    \item \texttt{PASS} : La directive \texttt{aaa local authentication attempts max-fail} est présente avec une valeur $\leq 5$.
    \item \texttt{FAIL} : Commande absente ou valeur supérieure à 5.
  \end{itemize}

  \textbf{Sévérité :} \textcolor{accentOrange}{\textbf{ÉLEVÉE}}
\end{rulebox}

\subsubsection{No Service Password-Recovery}

\begin{rulebox}{CISCO-MGT-PWD-003 — No Service Password-Recovery}
  \textbf{Objectif :} Désactiver la procédure de récupération de mot de passe par accès physique (console) pour les environnements à haute sécurité.

  \textbf{Commandes à vérifier :}
  \begin{lstlisting}
no service password-recovery
  \end{lstlisting}

  \textbf{Sévérité :} \textcolor{accentOrange}{\textbf{ÉLEVÉE}} — \textit{Applicable uniquement si la sécurité physique est garantie.}
\end{rulebox}

\subsection{Disable Unused Services}

\begin{rulebox}{CISCO-MGT-SVC-001 — Disable Unused Services}
  \textbf{Objectif :} Désactiver tous les services inutilisés pour réduire la surface d'attaque.

  \textbf{Services à désactiver (vérification de leur absence ou désactivation explicite) :}
  \begin{lstlisting}
no service tcp-small-servers
no service udp-small-servers
no service finger
no ip http server           ! Sauf si HTTPS gestion requise
no ip bootp server
no cdp run                  ! Si CDP non necessaire
no ip proxy-arp
no ip source-route
  \end{lstlisting}

  \textbf{Sévérité :} \textcolor{accentOrange}{\textbf{ÉLEVÉE}}
\end{rulebox}

\subsection{EXEC Timeout}

\begin{rulebox}{CISCO-MGT-EXEC-001 — EXEC Timeout}
  \textbf{Objectif :} Configurer un délai d'expiration des sessions EXEC pour éviter qu'une session inactive reste ouverte indéfiniment.

  \textbf{Commandes à vérifier (sur chaque ligne vty et console) :}
  \begin{lstlisting}
line con 0
 exec-timeout 5 0

line vty 0 4
 exec-timeout 10 0
  \end{lstlisting}

  \textbf{Logique :}
  \begin{itemize}
    \item \texttt{PASS} : \texttt{exec-timeout} présent sur toutes les lignes vty et console, valeur $\leq 10$ minutes.
    \item \texttt{FAIL} : Valeur = 0 0 (jamais de timeout) ou commande absente.
  \end{itemize}

  \textbf{Sévérité :} \textcolor{accentOrange}{\textbf{ÉLEVÉE}}
\end{rulebox}

\subsection{Keepalives for TCP Sessions}

\begin{rulebox}{CISCO-MGT-KEEPALIVE-001 — Keepalives for TCP Sessions}
  \textbf{Objectif :} Activer les keepalives TCP pour détecter les sessions TCP orphelines et libérer les ressources.

  \textbf{Commandes à vérifier :}
  \begin{lstlisting}
service tcp-keepalives-in
service tcp-keepalives-out
  \end{lstlisting}

  \textbf{Sévérité :} \textcolor{accentGreen}{\textbf{MOYENNE}}
\end{rulebox}

\subsection{Management Interface Use}

\begin{rulebox}{CISCO-MGT-INTF-001 — Management Interface Use}
  \textbf{Objectif :} Isoler le trafic de management sur une interface dédiée (VRF management) pour limiter l'exposition.

  \textbf{Commandes à vérifier :}
  \begin{lstlisting}
interface GigabitEthernetX
 vrf forwarding Mgmt-intf
 ip address <mgmt-ip> <mask>
  \end{lstlisting}

  \textbf{Sévérité :} \textcolor{accentOrange}{\textbf{ÉLEVÉE}}
\end{rulebox}

\subsection{Memory Threshold Notifications}

\begin{rulebox}{CISCO-MGT-MEM-001 — Memory Threshold Notifications}
  \textbf{Objectif :} Configurer des alertes lorsque la mémoire disponible descend en dessous d'un seuil critique.

  \textbf{Commandes à vérifier :}
  \begin{lstlisting}
memory free low-watermark processor <bytes>
memory free low-watermark io <bytes>
  \end{lstlisting}

  \textbf{Sévérité :} \textcolor{accentGreen}{\textbf{MOYENNE}}
\end{rulebox}

\subsection{CPU Thresholding Notification}

\begin{rulebox}{CISCO-MGT-CPU-001 — CPU Thresholding Notification}
  \textbf{Objectif :} Alerter lorsque l'utilisation CPU dépasse un seuil défini, pouvant indiquer une attaque DoS.

  \textbf{Commandes à vérifier :}
  \begin{lstlisting}
process cpu threshold type total rising <pct> interval <sec>
  \end{lstlisting}

  \textbf{Sévérité :} \textcolor{accentGreen}{\textbf{MOYENNE}}
\end{rulebox}

\subsection{Reserve Memory for Console Access}

\begin{rulebox}{CISCO-MGT-RES-001 — Reserve Memory for Console Access}
  \textbf{Objectif :} Réserver une portion de mémoire pour garantir l'accès console même en situation de pénurie mémoire.

  \begin{lstlisting}
memory reserve critical <bytes>
  \end{lstlisting}

  \textbf{Sévérité :} \textcolor{accentGreen}{\textbf{MOYENNE}}
\end{rulebox}

\subsection{Network Time Protocol (NTP)}

\begin{rulebox}{CISCO-MGT-NTP-001 — Network Time Protocol}
  \textbf{Objectif :} Synchroniser l'horloge de l'équipement avec un serveur NTP fiable et authentifié pour que les logs soient horodatés correctement.

  \textbf{Commandes à vérifier :}
  \begin{lstlisting}
ntp authenticate
ntp authentication-key <keyid> md5 <key>
ntp trusted-key <keyid>
ntp server <ip> key <keyid>
  \end{lstlisting}

  \textbf{Sévérité :} \textcolor{accentOrange}{\textbf{ÉLEVÉE}}
\end{rulebox}

\subsection{Disable Smart Install}

\begin{rulebox}{CISCO-MGT-SMI-001 — Disable Smart Install}
  \textbf{Objectif :} Désactiver le service Smart Install (SMI) s'il n'est pas utilisé, car il a été exploité massivement (CVE-2018-0171).

  \textbf{Commandes à vérifier :}
  \begin{lstlisting}
no vstack
  \end{lstlisting}

  \textbf{Logique :}
  \begin{itemize}
    \item \texttt{PASS} : \texttt{no vstack} présent dans le running-config.
    \item \texttt{FAIL} : \texttt{vstack} présent ou \texttt{no vstack} absent.
  \end{itemize}

  \textbf{Sévérité :} \textcolor{accentRed}{\textbf{CRITIQUE}}
\end{rulebox}

% ---- 3.2 ----
\section{Limit Access to the Network with Infrastructure ACLs}

Les Infrastructure ACLs (iACLs) filtrent le trafic destiné aux adresses IP de l'infrastructure (interfaces de management, loopbacks) afin de n'autoriser que les sources légitimes.

\subsection{ICMP Packet Filtering}

\begin{rulebox}{CISCO-ACL-ICMP-001 — ICMP Packet Filtering}
  \textbf{Objectif :} Restreindre les paquets ICMP entrants vers les adresses de management aux seuls hôtes autorisés.

  \textbf{Extrait d'iACL attendu :}
  \begin{lstlisting}
ip access-list extended iACL-INBOUND
 permit icmp <trusted-hosts> host <router-ip> echo
 permit icmp <trusted-hosts> host <router-ip> echo-reply
 deny   icmp any any log
  \end{lstlisting}

  \textbf{Sévérité :} \textcolor{accentOrange}{\textbf{ÉLEVÉE}}
\end{rulebox}

\subsection{Filter IP Fragments}

\begin{rulebox}{CISCO-ACL-FRAG-001 — Filter IP Fragments}
  \textbf{Objectif :} Bloquer les fragments IP non initiaux qui peuvent être utilisés pour contourner les ACL et mener des attaques.

  \begin{lstlisting}
deny ip any any fragments log
  \end{lstlisting}

  \textbf{Sévérité :} \textcolor{accentOrange}{\textbf{ÉLEVÉE}}
\end{rulebox}

\subsection{ACL Support to Filter IP Options}

\begin{rulebox}{CISCO-ACL-IPOPT-001 — ACL Support to Filter IP Options}
  \textbf{Objectif :} Bloquer les paquets IP avec des options inhabituelles (Record Route, Timestamp, Loose/Strict Source Routing).

  \begin{lstlisting}
ip access-list extended iACL-INBOUND
 deny ip any any option record-route log
 deny ip any any option timestamp  log
 deny ip any any option traceroute  log
  \end{lstlisting}

  \textbf{Sévérité :} \textcolor{accentOrange}{\textbf{ÉLEVÉE}}
\end{rulebox}

\subsection{ACL Support to Filter on TTL Value}

\begin{rulebox}{CISCO-ACL-TTL-001 — Filter on TTL Value}
  \textbf{Objectif :} Rejeter les paquets avec TTL=0 ou TTL=1 non légitimes pour éviter certaines reconnaissances réseau.

  \begin{lstlisting}
ip access-list extended iACL-INBOUND
 deny ip any any ttl lt 2 log
  \end{lstlisting}

  \textbf{Sévérité :} \textcolor{accentGreen}{\textbf{MOYENNE}}
\end{rulebox}

% ---- 3.3 ----
\section{Secure Interactive Management Sessions}

\subsection{Management Plane Protection}

\begin{rulebox}{CISCO-MGT-MPP-001 — Management Plane Protection}
  \textbf{Objectif :} Restreindre les protocoles de management à une interface spécifique (interface de management dédiée).

  \begin{lstlisting}
control-plane host
 management-interface GigabitEthernetX allow ssh https
  \end{lstlisting}

  \textbf{Sévérité :} \textcolor{accentRed}{\textbf{CRITIQUE}}
\end{rulebox}

\subsection{SSHv2}

\begin{rulebox}{CISCO-MGT-SSH-001 — SSHv2 Activation}
  \textbf{Objectif :} S'assurer que seul SSHv2 est activé pour l'administration distante. Telnet doit être désactivé.

  \textbf{Commandes à vérifier :}
  \begin{lstlisting}
ip ssh version 2
ip ssh time-out 60
ip ssh authentication-retries 3

line vty 0 4
 transport input ssh    ! Pas de telnet
  \end{lstlisting}

  \textbf{Logique :}
  \begin{itemize}
    \item \texttt{PASS} : \texttt{ip ssh version 2} présent ET \texttt{transport input ssh} (sans telnet) sur toutes les vty.
    \item \texttt{FAIL} : \texttt{transport input telnet} détecté OU \texttt{ip ssh version 2} absent.
  \end{itemize}

  \textbf{Sévérité :} \textcolor{accentRed}{\textbf{CRITIQUE}}
\end{rulebox}

\subsection{SSHv2 Enhancements for RSA Keys}

\begin{rulebox}{CISCO-MGT-SSH-002 — RSA Key Size}
  \textbf{Objectif :} S'assurer que les clés RSA utilisées pour SSH ont une taille minimale de 2048 bits.

  \begin{lstlisting}
crypto key generate rsa modulus 2048
ip ssh rsa keypair-name <name>
  \end{lstlisting}

  \textbf{Sévérité :} \textcolor{accentOrange}{\textbf{ÉLEVÉE}}
\end{rulebox}

\subsection{Console and AUX Ports}

\begin{rulebox}{CISCO-MGT-CON-001 — Console and AUX Ports}
  \textbf{Objectif :} Sécuriser ou désactiver les ports Console et AUX.

  \begin{lstlisting}
line con 0
 exec-timeout 5 0
 login authentication <list>
 transport output ssh

line aux 0
 exec-timeout 0 1
 no exec
 transport input none
  \end{lstlisting}

  \textbf{Sévérité :} \textcolor{accentOrange}{\textbf{ÉLEVÉE}}
\end{rulebox}

\subsection{Control vty and tty Lines}

\begin{rulebox}{CISCO-MGT-VTY-001 — Control vty and tty Lines}
  \textbf{Objectif :} Limiter le nombre de lignes vty actives et appliquer des ACL d'accès.

  \begin{lstlisting}
line vty 0 4
 access-class <ACL-MGMT> in
 login authentication default
 exec-timeout 10 0
  \end{lstlisting}

  \textbf{Sévérité :} \textcolor{accentRed}{\textbf{CRITIQUE}}
\end{rulebox}

\subsection{Warning Banners}

\begin{rulebox}{CISCO-MGT-BANNER-001 — Warning Banners}
  \textbf{Objectif :} Afficher une bannière d'avertissement légale avant toute connexion.

  \begin{lstlisting}
banner motd ^
  AVERTISSEMENT : Acces autorise uniquement.
  Toute connexion non autorisee est interdite
  et sera poursuivie.
^
  \end{lstlisting}

  \textbf{Sévérité :} \textcolor{accentGreen}{\textbf{MOYENNE}}
\end{rulebox}

% ---- 3.4 ----
\section{Authentication, Authorization, and Accounting (AAA)}

\subsection{TACACS+ Authentication}

\begin{rulebox}{CISCO-AAA-TACACS-001 — TACACS+ Authentication}
  \textbf{Objectif :} Centraliser l'authentification via TACACS+ permettant un contrôle granulaire et une traçabilité complète.

  \begin{lstlisting}
aaa new-model
tacacs server <name>
 address ipv4 <tacacs-server-ip>
 key <secret>
aaa authentication login default group tacacs+ local
  \end{lstlisting}

  \textbf{Sévérité :} \textcolor{accentRed}{\textbf{CRITIQUE}}
\end{rulebox}

\subsection{Authentication Fallback}

\begin{rulebox}{CISCO-AAA-FALLBACK-001 — Authentication Fallback}
  \textbf{Objectif :} Configurer un compte local de secours en cas d'indisponibilité du serveur TACACS+.

  \begin{lstlisting}
aaa authentication login default group tacacs+ local
username <admin> privilege 15 secret <strong-password>
  \end{lstlisting}

  \textbf{Sévérité :} \textcolor{accentOrange}{\textbf{ÉLEVÉE}}
\end{rulebox}

\subsection{Use of Type 7 Passwords}

\begin{rulebox}{CISCO-AAA-TYPE7-001 — Type 7 Password Detection}
  \textbf{Objectif :} Détecter l'utilisation de mots de passe de type 7 (facilement déchiffrables) et les signaler.

  \textbf{Logique de détection :}
  \begin{itemize}
    \item Rechercher les patterns \texttt{password 7 <hash>} dans le running-config.
    \item \texttt{FAIL} si des mots de passe type 7 sont trouvés pour des lignes vty/console ou des comptes utilisateurs.
  \end{itemize}

  \textbf{Sévérité :} \textcolor{accentRed}{\textbf{CRITIQUE}}
\end{rulebox}

\subsection{TACACS+ Command Authorization}

\begin{rulebox}{CISCO-AAA-AUTHZ-001 — TACACS+ Command Authorization}
  \textbf{Objectif :} Activer l'autorisation de commandes via TACACS+ pour contrôler quelles commandes chaque utilisateur peut exécuter.

  \begin{lstlisting}
aaa authorization commands 1 default group tacacs+ local
aaa authorization commands 15 default group tacacs+ local
  \end{lstlisting}

  \textbf{Sévérité :} \textcolor{accentOrange}{\textbf{ÉLEVÉE}}
\end{rulebox}

\subsection{TACACS+ Command Accounting}

\begin{rulebox}{CISCO-AAA-ACCT-001 — TACACS+ Command Accounting}
  \textbf{Objectif :} Enregistrer toutes les commandes exécutées sur le serveur TACACS+ à des fins d'audit.

  \begin{lstlisting}
aaa accounting commands 1 default start-stop group tacacs+
aaa accounting commands 15 default start-stop group tacacs+
aaa accounting exec default start-stop group tacacs+
  \end{lstlisting}

  \textbf{Sévérité :} \textcolor{accentOrange}{\textbf{ÉLEVÉE}}
\end{rulebox}

\subsection{Redundant AAA Servers}

\begin{rulebox}{CISCO-AAA-REDUND-001 — Redundant AAA Servers}
  \textbf{Objectif :} Configurer au moins deux serveurs AAA pour assurer la haute disponibilité.

  \begin{lstlisting}
aaa group server tacacs+ TACACS-SERVERS
 server name TACACS-PRIMARY
 server name TACACS-SECONDARY
  \end{lstlisting}

  \textbf{Sévérité :} \textcolor{accentGreen}{\textbf{MOYENNE}}
\end{rulebox}

% ---- 3.5 ----
\section{Fortify the Simple Network Management Protocol (SNMP)}

\subsection{SNMP Community Strings}

\begin{rulebox}{CISCO-SNMP-COMM-001 — SNMP Community Strings}
  \textbf{Objectif :} S'assurer que les community strings SNMP par défaut (\texttt{public}, \texttt{private}) ne sont pas utilisées.

  \textbf{Logique :}
  \begin{itemize}
    \item \texttt{FAIL} si \texttt{snmp-server community public} ou \texttt{snmp-server community private} détecté.
  \end{itemize}

  \textbf{Sévérité :} \textcolor{accentRed}{\textbf{CRITIQUE}}
\end{rulebox}

\subsection{SNMP Community Strings with ACLs}

\begin{rulebox}{CISCO-SNMP-COMM-002 — SNMP Community Strings with ACLs}
  \textbf{Objectif :} Associer une ACL à chaque community string SNMP pour restreindre l'accès.

  \begin{lstlisting}
snmp-server community <string> RO <ACL-NUMBER>
  \end{lstlisting}

  \textbf{Sévérité :} \textcolor{accentOrange}{\textbf{ÉLEVÉE}}
\end{rulebox}

\subsection{SNMP Version 3}

\begin{rulebox}{CISCO-SNMP-V3-001 — SNMP Version 3}
  \textbf{Objectif :} Utiliser SNMPv3 avec authentification et chiffrement (authPriv) plutôt que SNMPv1/v2c.

  \begin{lstlisting}
snmp-server group <group> v3 priv
snmp-server user <user> <group> v3 auth sha <authpass> priv aes 128 <privpass>
  \end{lstlisting}

  \textbf{Sévérité :} \textcolor{accentRed}{\textbf{CRITIQUE}}
\end{rulebox}

% ---- 3.6 ----
\section{Logging Best Practices}

\subsection{Send Logs to a Central Location}

\begin{rulebox}{CISCO-LOG-SYSLOG-001 — Send Logs to a Central Location}
  \textbf{Objectif :} Envoyer les logs vers un serveur Syslog centralisé.

  \begin{lstlisting}
logging host <syslog-server-ip>
logging trap informational
  \end{lstlisting}

  \textbf{Sévérité :} \textcolor{accentOrange}{\textbf{ÉLEVÉE}}
\end{rulebox}

\subsection{Logging Level}

\begin{rulebox}{CISCO-LOG-LEVEL-001 — Logging Level}
  \textbf{Objectif :} Configurer un niveau de logging approprié (informational ou debugging en production selon besoin).

  \begin{lstlisting}
logging trap informational   ! Niveau 6 minimum
  \end{lstlisting}

  \textbf{Sévérité :} \textcolor{accentGreen}{\textbf{MOYENNE}}
\end{rulebox}

\subsection{Do Not Log to Console or Monitor Sessions}

\begin{rulebox}{CISCO-LOG-CONSOLE-001 — No Console/Monitor Logging}
  \textbf{Objectif :} Désactiver le logging vers la console et les sessions monitor pour éviter la dégradation des performances.

  \begin{lstlisting}
no logging console
no logging monitor
  \end{lstlisting}

  \textbf{Sévérité :} \textcolor{accentGreen}{\textbf{MOYENNE}}
\end{rulebox}

\subsection{Use Buffered Logs}

\begin{rulebox}{CISCO-LOG-BUFFER-001 — Use Buffered Logs}
  \textbf{Objectif :} Maintenir un buffer de logs local suffisant.

  \begin{lstlisting}
logging buffered 64000 informational
  \end{lstlisting}

  \textbf{Sévérité :} \textcolor{accentGreen}{\textbf{MOYENNE}}
\end{rulebox}

\subsection{Configure Logging Source Interface}

\begin{rulebox}{CISCO-LOG-SRCINTF-001 — Configure Logging Source Interface}
  \textbf{Objectif :} Utiliser une interface source stable (loopback) pour les messages syslog afin de garantir leur traçabilité.

  \begin{lstlisting}
logging source-interface Loopback0
  \end{lstlisting}

  \textbf{Sévérité :} \textcolor{accentGreen}{\textbf{MOYENNE}}
\end{rulebox}

\subsection{Configure Log Timestamps}

\begin{rulebox}{CISCO-LOG-TIMESTAMP-001 — Configure Log Timestamps}
  \textbf{Objectif :} Horodater les logs avec une précision à la milliseconde et inclure le fuseau horaire.

  \begin{lstlisting}
service timestamps log datetime msec localtime show-timezone
service timestamps debug datetime msec localtime show-timezone
  \end{lstlisting}

  \textbf{Sévérité :} \textcolor{accentOrange}{\textbf{ÉLEVÉE}}
\end{rulebox}

% ---- 3.7 ----
\section{Cisco IOS Software Configuration Management}

\subsection{Configuration Replace and Rollback}

\begin{rulebox}{CISCO-CFG-ROLLBACK-001 — Configuration Replace and Rollback}
  \textbf{Objectif :} Activer la fonctionnalité de remplacement et de rollback de configuration.

  \begin{lstlisting}
archive
 path flash:archive
 maximum 14
 time-period 1440
  \end{lstlisting}

  \textbf{Sévérité :} \textcolor{accentGreen}{\textbf{MOYENNE}}
\end{rulebox}

\subsection{Cisco IOS Software Resilient Configuration}

\begin{rulebox}{CISCO-CFG-RESILIENT-001 — Resilient Configuration}
  \textbf{Objectif :} Activer la configuration résiliente pour protéger les fichiers IOS et running-config contre la suppression.

  \begin{lstlisting}
secure boot-image
secure boot-config
  \end{lstlisting}

  \textbf{Sévérité :} \textcolor{accentOrange}{\textbf{ÉLEVÉE}}
\end{rulebox}

% ============================================================
\chapter{Control Plane — Règles de Conformité}
% ============================================================

Le \textbf{Control Plane} traite les paquets nécessaires au fonctionnement des protocoles réseau (routage, spanning tree, etc.). Sa protection est essentielle pour la stabilité du réseau.

\section{General Control Plane Hardening}

\subsection{IP ICMP Redirects}

\begin{rulebox}{CISCO-CP-ICMP-001 — Disable IP ICMP Redirects}
  \textbf{Objectif :} Désactiver les ICMP Redirects sur toutes les interfaces pour éviter que l'équipement informe les hôtes de meilleures routes (exploitable).

  \begin{lstlisting}
interface <all>
 no ip redirects
  \end{lstlisting}

  \textbf{Logique :} Vérifier l'absence de \texttt{ip redirects} (ou présence de \texttt{no ip redirects}) sur toutes les interfaces.

  \textbf{Sévérité :} \textcolor{accentOrange}{\textbf{ÉLEVÉE}}
\end{rulebox}

\subsection{ICMP Unreachables}

\begin{rulebox}{CISCO-CP-ICMP-002 — Disable ICMP Unreachables}
  \textbf{Objectif :} Désactiver les messages ICMP Unreachable sur les interfaces exposées pour limiter les informations données aux attaquants.

  \begin{lstlisting}
interface <external>
 no ip unreachables
  \end{lstlisting}

  \textbf{Sévérité :} \textcolor{accentGreen}{\textbf{MOYENNE}}
\end{rulebox}

\subsection{Proxy ARP}

\begin{rulebox}{CISCO-CP-PROXYARP-001 — Disable Proxy ARP}
  \textbf{Objectif :} Désactiver le Proxy ARP pour éviter que le routeur réponde aux requêtes ARP à la place d'autres hôtes.

  \begin{lstlisting}
interface <all>
 no ip proxy-arp
  \end{lstlisting}

  \textbf{Sévérité :} \textcolor{accentGreen}{\textbf{MOYENNE}}
\end{rulebox}

\section{Limit CPU Impact of Control Plane Traffic}

\subsection{Control Plane Policing (CoPP)}

\begin{rulebox}{CISCO-CP-COPP-001 — Control Plane Policing}
  \textbf{Objectif :} Implémenter une politique CoPP pour protéger le CPU du routeur contre les inondations de paquets.

  \begin{lstlisting}
policy-map COPP-POLICY
 class CRITICAL-TRAFFIC
  police rate 32000 bps conform-action transmit
            exceed-action drop

control-plane
 service-policy input COPP-POLICY
  \end{lstlisting}

  \textbf{Logique :} Vérifier la présence d'une \texttt{service-policy input} sur le \texttt{control-plane}.

  \textbf{Sévérité :} \textcolor{accentRed}{\textbf{CRITIQUE}}
\end{rulebox}

\subsection{Receive ACLs (rACL)}

\begin{rulebox}{CISCO-CP-RACL-001 — Receive ACLs}
  \textbf{Objectif :} Utiliser des Receive ACLs (disponibles sur certaines plateformes) pour filtrer le trafic destiné au routeur avant traitement CPU.

  \textbf{Sévérité :} \textcolor{accentOrange}{\textbf{ÉLEVÉE}}
\end{rulebox}

\section{Secure BGP}

\subsection{TTL-based Security Protections (GTSM)}

\begin{rulebox}{CISCO-BGP-GTSM-001 — TTL-based Security (GTSM)}
  \textbf{Objectif :} Utiliser GTSM (Generalized TTL Security Mechanism) pour n'accepter les sessions BGP que de voisins à 1 saut.

  \begin{lstlisting}
router bgp <AS>
 neighbor <peer-ip> ttl-security hops 1
  \end{lstlisting}

  \textbf{Sévérité :} \textcolor{accentOrange}{\textbf{ÉLEVÉE}}
\end{rulebox}

\subsection{BGP Peer Authentication with MD5}

\begin{rulebox}{CISCO-BGP-MD5-001 — BGP Peer Authentication}
  \textbf{Objectif :} Authentifier toutes les sessions BGP par mot de passe MD5 pour éviter les injections de routes.

  \begin{lstlisting}
router bgp <AS>
 neighbor <peer-ip> password <secret>
  \end{lstlisting}

  \textbf{Sévérité :} \textcolor{accentOrange}{\textbf{ÉLEVÉE}}
\end{rulebox}

\subsection{Configure Maximum Prefixes}

\begin{rulebox}{CISCO-BGP-MAXPFX-001 — Maximum Prefixes}
  \textbf{Objectif :} Limiter le nombre de préfixes reçus d'un pair BGP pour éviter les tables de routage explosives.

  \begin{lstlisting}
router bgp <AS>
 neighbor <peer-ip> maximum-prefix <N> <warning-pct>
  \end{lstlisting}

  \textbf{Sévérité :} \textcolor{accentOrange}{\textbf{ÉLEVÉE}}
\end{rulebox}

\subsection{Filter BGP Prefixes with Prefix Lists}

\begin{rulebox}{CISCO-BGP-PFXLIST-001 — Prefix List Filtering}
  \textbf{Objectif :} Filtrer les préfixes BGP entrants et sortants via des prefix-lists pour n'accepter que les routes légitimes.

  \begin{lstlisting}
router bgp <AS>
 neighbor <peer-ip> prefix-list IN-FILTER in
 neighbor <peer-ip> prefix-list OUT-FILTER out
  \end{lstlisting}

  \textbf{Sévérité :} \textcolor{accentOrange}{\textbf{ÉLEVÉE}}
\end{rulebox}

\section{Secure Interior Gateway Protocols}

\subsection{Routing Protocol Authentication with MD5}

\begin{rulebox}{CISCO-IGP-MD5-001 — IGP Authentication}
  \textbf{Objectif :} Authentifier les échanges OSPF, EIGRP, IS-IS par clé MD5 ou SHA pour éviter les injections de routes.

  \textbf{OSPF :}
  \begin{lstlisting}
router ospf <process>
 area <X> authentication message-digest
interface <X>
 ip ospf message-digest-key 1 md5 <key>
  \end{lstlisting}

  \textbf{EIGRP :}
  \begin{lstlisting}
key chain EIGRP-KEYS
 key 1
  key-string <secret>
interface <X>
 ip authentication mode eigrp <AS> md5
 ip authentication key-chain eigrp <AS> EIGRP-KEYS
  \end{lstlisting}

  \textbf{Sévérité :} \textcolor{accentRed}{\textbf{CRITIQUE}}
\end{rulebox}

\subsection{Passive-Interface Commands}

\begin{rulebox}{CISCO-IGP-PASSIVE-001 — Passive-Interface}
  \textbf{Objectif :} Configurer les interfaces non routables en mode passif pour éviter l'envoi de mises à jour de routage sur ces segments.

  \begin{lstlisting}
router ospf <process>
 passive-interface default
 no passive-interface <routing-interface>
  \end{lstlisting}

  \textbf{Sévérité :} \textcolor{accentGreen}{\textbf{MOYENNE}}
\end{rulebox}

\section{Secure First Hop Redundancy Protocols (FHRP)}

\begin{rulebox}{CISCO-FHRP-AUTH-001 — FHRP Authentication}
  \textbf{Objectif :} Authentifier les messages HSRP, VRRP, GLBP pour éviter qu'un attaquant devienne le routeur actif.

  \textbf{HSRP :}
  \begin{lstlisting}
interface <X>
 standby <group> authentication md5 key-string <secret>
  \end{lstlisting}

  \textbf{Sévérité :} \textcolor{accentOrange}{\textbf{ÉLEVÉE}}
\end{rulebox}

% ============================================================
\chapter{Data Plane — Règles de Conformité}
% ============================================================

Le \textbf{Data Plane} traite le trafic de transit. Sa sécurisation vise à prévenir le spoofing, les attaques DoS et les abus des ressources du plan de données.

\section{General Data Plane Hardening}

\subsection{IP Options Selective Drop}

\begin{rulebox}{CISCO-DP-IPOPT-001 — IP Options Selective Drop}
  \textbf{Objectif :} Activer le filtrage sélectif des paquets IP avec des options pour réduire la charge CPU.

  \begin{lstlisting}
ip options drop
  \end{lstlisting}

  \textbf{Sévérité :} \textcolor{accentGreen}{\textbf{MOYENNE}}
\end{rulebox}

\subsection{Disable IP Source Routing}

\begin{rulebox}{CISCO-DP-SRCROUTE-001 — Disable IP Source Routing}
  \textbf{Objectif :} Désactiver le Source Routing IP (Loose et Strict) pour empêcher qu'un attaquant dicte le chemin de ses paquets.

  \begin{lstlisting}
no ip source-route
  \end{lstlisting}

  \textbf{Sévérité :} \textcolor{accentOrange}{\textbf{ÉLEVÉE}}
\end{rulebox}

\subsection{Disable ICMP Redirects}

\begin{rulebox}{CISCO-DP-ICMPRED-001 — Disable ICMP Redirects (Data Plane)}
  \textbf{Objectif :} Désactiver les ICMP Redirects sur toutes les interfaces (complément du plan de contrôle).

  \begin{lstlisting}
interface <all>
 no ip redirects
  \end{lstlisting}

  \textbf{Sévérité :} \textcolor{accentOrange}{\textbf{ÉLEVÉE}}
\end{rulebox}

\subsection{Disable or Limit IP Directed Broadcasts}

\begin{rulebox}{CISCO-DP-DIRCAST-001 — Disable IP Directed Broadcasts}
  \textbf{Objectif :} Désactiver les broadcasts dirigés IP pour prévenir les attaques Smurf.

  \begin{lstlisting}
interface <all>
 no ip directed-broadcast
  \end{lstlisting}

  \textbf{Sévérité :} \textcolor{accentOrange}{\textbf{ÉLEVÉE}}
\end{rulebox}

\section{Filter Transit Traffic with Transit ACLs}

Les Transit ACLs (tACL) filtrent le trafic de transit traversant l'équipement vers les segments internes critiques.

\subsection{ICMP Packet Filtering (Transit)}

\begin{rulebox}{CISCO-TACL-ICMP-001 — Transit ICMP Filtering}
  \textbf{Objectif :} Limiter les types ICMP autorisés en transit.

  \begin{lstlisting}
ip access-list extended tACL
 permit icmp any any echo
 permit icmp any any echo-reply
 permit icmp any any unreachable
 deny   icmp any any log
  \end{lstlisting}

  \textbf{Sévérité :} \textcolor{accentOrange}{\textbf{ÉLEVÉE}}
\end{rulebox}

\section{Anti-Spoofing Protections}

\subsection{Unicast Reverse Path Forwarding (uRPF)}

\begin{rulebox}{CISCO-ANTISPOOFING-URPF-001 — Unicast RPF}
  \textbf{Objectif :} Activer uRPF sur les interfaces externes pour rejeter les paquets dont l'adresse source ne correspond pas à la table de routage.

  \begin{lstlisting}
interface <external>
 ip verify unicast source reachable-via rx
  \end{lstlisting}

  \textbf{Sévérité :} \textcolor{accentRed}{\textbf{CRITIQUE}}
\end{rulebox}

\subsection{IP Source Guard}

\begin{rulebox}{CISCO-ANTISPOOFING-IPSG-001 — IP Source Guard}
  \textbf{Objectif :} Activer IP Source Guard sur les ports d'accès pour n'autoriser que les paquets correspondant aux liaisons DHCP.

  \begin{lstlisting}
interface <access-port>
 ip verify source
  \end{lstlisting}

  \textbf{Sévérité :} \textcolor{accentOrange}{\textbf{ÉLEVÉE}}
\end{rulebox}

\subsection{Dynamic ARP Inspection (DAI)}

\begin{rulebox}{CISCO-ANTISPOOFING-DAI-001 — Dynamic ARP Inspection}
  \textbf{Objectif :} Activer DAI sur les VLANs pour prévenir le poisoning ARP.

  \begin{lstlisting}
ip arp inspection vlan <vlan-id>
interface <trunk>
 ip arp inspection trust
  \end{lstlisting}

  \textbf{Sévérité :} \textcolor{accentOrange}{\textbf{ÉLEVÉE}}
\end{rulebox}

\subsection{Port Security}

\begin{rulebox}{CISCO-ANTISPOOFING-PORTSEC-001 — Port Security}
  \textbf{Objectif :} Limiter le nombre d'adresses MAC autorisées par port pour prévenir les attaques MAC flooding.

  \begin{lstlisting}
interface <access-port>
 switchport port-security maximum 2
 switchport port-security violation restrict
 switchport port-security
  \end{lstlisting}

  \textbf{Sévérité :} \textcolor{accentGreen}{\textbf{MOYENNE}}
\end{rulebox}

\section{Traffic Identification and Traceback}

\subsection{NetFlow}

\begin{rulebox}{CISCO-TRAFFIC-NETFLOW-001 — NetFlow}
  \textbf{Objectif :} Activer NetFlow pour la visibilité du trafic et la détection d'anomalies.

  \begin{lstlisting}
flow record NETFLOW-RECORD
 match ipv4 source address
 match ipv4 destination address
 match transport source-port
 collect counter bytes

flow exporter NETFLOW-EXPORT
 destination <collector-ip>
 transport udp 2055
  \end{lstlisting}

  \textbf{Sévérité :} \textcolor{accentGreen}{\textbf{MOYENNE}}
\end{rulebox}

\section{Access Control with VLAN Maps and PACLs}

\subsection{Access Control with VLAN Maps}

\begin{rulebox}{CISCO-VLAN-MAP-001 — VLAN Maps Access Control}
  \textbf{Objectif :} Utiliser des VLAN Access Maps pour contrôler le trafic intra-VLAN.

  \begin{lstlisting}
vlan access-map VLAN-FILTER 10
 match ip address <ACL>
 action forward

vlan filter VLAN-FILTER vlan-list <vlan-id>
  \end{lstlisting}

  \textbf{Sévérité :} \textcolor{accentGreen}{\textbf{MOYENNE}}
\end{rulebox}

\section{Private VLAN Use}

\subsection{Isolated VLANs}

\begin{rulebox}{CISCO-PVLAN-ISO-001 — Isolated VLANs}
  \textbf{Objectif :} Utiliser les VLANs privés isolés pour empêcher la communication directe entre hôtes d'un même VLAN.

  \begin{lstlisting}
vlan <isolated-vlan-id>
 private-vlan isolated

vlan <primary-vlan-id>
 private-vlan primary
 private-vlan association <isolated-vlan-id>
  \end{lstlisting}

  \textbf{Sévérité :} \textcolor{accentGreen}{\textbf{MOYENNE}}
\end{rulebox}

% ============================================================
\chapter{Tableau Récapitulatif des Règles}
% ============================================================

\begin{longtable}{>{\small}p{3.8cm} >{\small}p{3cm} >{\small}p{2cm} >{\small}p{2cm} >{\small}p{2.5cm}}
  \caption{Récapitulatif complet des règles de conformité Cisco} \\
  \toprule
  \rowcolor{tableHeader}
  \color{white}\textbf{ID Règle} &
  \color{white}\textbf{Titre} &
  \color{white}\textbf{Plan} &
  \color{white}\textbf{Sévérité} &
  \color{white}\textbf{Type de check} \\
  \midrule
  \endfirsthead
  \multicolumn{5}{c}{\small\textit{(suite du tableau)}} \\
  \toprule
  \rowcolor{tableHeader}
  \color{white}\textbf{ID Règle} &
  \color{white}\textbf{Titre} &
  \color{white}\textbf{Plan} &
  \color{white}\textbf{Sévérité} &
  \color{white}\textbf{Type de check} \\
  \midrule
  \endhead
  \bottomrule
  \endfoot

  \rowcolor{tableRowA} CISCO-MGT-PWD-001 & Enhanced Password Security & Management & \textcolor{accentRed}{CRITIQUE} & regex\_present \\
  \rowcolor{tableRowB} CISCO-MGT-PWD-002 & Login Retry Lockout & Management & \textcolor{accentOrange}{ÉLEVÉE} & regex\_value \\
  \rowcolor{tableRowA} CISCO-MGT-PWD-003 & No Service Password-Recovery & Management & \textcolor{accentOrange}{ÉLEVÉE} & regex\_present \\
  \rowcolor{tableRowB} CISCO-MGT-SVC-001 & Disable Unused Services & Management & \textcolor{accentOrange}{ÉLEVÉE} & multi\_regex \\
  \rowcolor{tableRowA} CISCO-MGT-EXEC-001 & EXEC Timeout & Management & \textcolor{accentOrange}{ÉLEVÉE} & regex\_value \\
  \rowcolor{tableRowB} CISCO-MGT-KEEPALIVE-001 & Keepalives TCP & Management & \textcolor{accentGreen}{MOYENNE} & regex\_present \\
  \rowcolor{tableRowA} CISCO-MGT-INTF-001 & Management Interface & Management & \textcolor{accentOrange}{ÉLEVÉE} & regex\_present \\
  \rowcolor{tableRowB} CISCO-MGT-MEM-001 & Memory Threshold & Management & \textcolor{accentGreen}{MOYENNE} & regex\_present \\
  \rowcolor{tableRowA} CISCO-MGT-CPU-001 & CPU Thresholding & Management & \textcolor{accentGreen}{MOYENNE} & regex\_present \\
  \rowcolor{tableRowB} CISCO-MGT-RES-001 & Reserve Memory Console & Management & \textcolor{accentGreen}{MOYENNE} & regex\_present \\
  \rowcolor{tableRowA} CISCO-MGT-NTP-001 & NTP Authenticated & Management & \textcolor{accentOrange}{ÉLEVÉE} & multi\_regex \\
  \rowcolor{tableRowB} CISCO-MGT-SMI-001 & Disable Smart Install & Management & \textcolor{accentRed}{CRITIQUE} & regex\_present \\
  \rowcolor{tableRowA} CISCO-ACL-ICMP-001 & ICMP Packet Filtering & Management & \textcolor{accentOrange}{ÉLEVÉE} & acl\_check \\
  \rowcolor{tableRowB} CISCO-ACL-FRAG-001 & Filter IP Fragments & Management & \textcolor{accentOrange}{ÉLEVÉE} & acl\_check \\
  \rowcolor{tableRowA} CISCO-ACL-IPOPT-001 & Filter IP Options & Management & \textcolor{accentOrange}{ÉLEVÉE} & acl\_check \\
  \rowcolor{tableRowB} CISCO-ACL-TTL-001 & Filter TTL Value & Management & \textcolor{accentGreen}{MOYENNE} & acl\_check \\
  \rowcolor{tableRowA} CISCO-MGT-MPP-001 & Management Plane Protection & Management & \textcolor{accentRed}{CRITIQUE} & regex\_present \\
  \rowcolor{tableRowB} CISCO-MGT-SSH-001 & SSHv2 Activation & Management & \textcolor{accentRed}{CRITIQUE} & multi\_regex \\
  \rowcolor{tableRowA} CISCO-MGT-SSH-002 & RSA Key Size 2048 & Management & \textcolor{accentOrange}{ÉLEVÉE} & regex\_value \\
  \rowcolor{tableRowB} CISCO-MGT-CON-001 & Console/AUX Ports & Management & \textcolor{accentOrange}{ÉLEVÉE} & multi\_regex \\
  \rowcolor{tableRowA} CISCO-MGT-VTY-001 & Control vty Lines & Management & \textcolor{accentRed}{CRITIQUE} & multi\_regex \\
  \rowcolor{tableRowB} CISCO-MGT-BANNER-001 & Warning Banners & Management & \textcolor{accentGreen}{MOYENNE} & regex\_present \\
  \rowcolor{tableRowA} CISCO-AAA-TACACS-001 & TACACS+ Authentication & Management & \textcolor{accentRed}{CRITIQUE} & multi\_regex \\
  \rowcolor{tableRowB} CISCO-AAA-FALLBACK-001 & Authentication Fallback & Management & \textcolor{accentOrange}{ÉLEVÉE} & regex\_present \\
  \rowcolor{tableRowA} CISCO-AAA-TYPE7-001 & No Type 7 Passwords & Management & \textcolor{accentRed}{CRITIQUE} & regex\_absent \\
  \rowcolor{tableRowB} CISCO-AAA-AUTHZ-001 & TACACS+ Command Authz & Management & \textcolor{accentOrange}{ÉLEVÉE} & regex\_present \\
  \rowcolor{tableRowA} CISCO-AAA-ACCT-001 & TACACS+ Accounting & Management & \textcolor{accentOrange}{ÉLEVÉE} & regex\_present \\
  \rowcolor{tableRowB} CISCO-AAA-REDUND-001 & Redundant AAA Servers & Management & \textcolor{accentGreen}{MOYENNE} & value\_count \\
  \rowcolor{tableRowA} CISCO-SNMP-COMM-001 & SNMP Default Strings & Management & \textcolor{accentRed}{CRITIQUE} & regex\_absent \\
  \rowcolor{tableRowB} CISCO-SNMP-COMM-002 & SNMP Community ACLs & Management & \textcolor{accentOrange}{ÉLEVÉE} & acl\_check \\
  \rowcolor{tableRowA} CISCO-SNMP-V3-001 & SNMPv3 authPriv & Management & \textcolor{accentRed}{CRITIQUE} & regex\_present \\
  \rowcolor{tableRowB} CISCO-LOG-SYSLOG-001 & Send Logs to Syslog & Management & \textcolor{accentOrange}{ÉLEVÉE} & regex\_present \\
  \rowcolor{tableRowA} CISCO-LOG-LEVEL-001 & Logging Level & Management & \textcolor{accentGreen}{MOYENNE} & regex\_value \\
  \rowcolor{tableRowB} CISCO-LOG-CONSOLE-001 & No Console Logging & Management & \textcolor{accentGreen}{MOYENNE} & regex\_present \\
  \rowcolor{tableRowA} CISCO-LOG-BUFFER-001 & Buffered Logs & Management & \textcolor{accentGreen}{MOYENNE} & regex\_present \\
  \rowcolor{tableRowB} CISCO-LOG-SRCINTF-001 & Logging Source Interface & Management & \textcolor{accentGreen}{MOYENNE} & regex\_present \\
  \rowcolor{tableRowA} CISCO-LOG-TIMESTAMP-001 & Log Timestamps & Management & \textcolor{accentOrange}{ÉLEVÉE} & regex\_present \\
  \rowcolor{tableRowB} CISCO-CFG-ROLLBACK-001 & Config Rollback & Management & \textcolor{accentGreen}{MOYENNE} & regex\_present \\
  \rowcolor{tableRowA} CISCO-CFG-RESILIENT-001 & Resilient Config & Management & \textcolor{accentOrange}{ÉLEVÉE} & regex\_present \\
  \rowcolor{tableRowB} CISCO-CP-ICMP-001 & Disable ICMP Redirects & Control & \textcolor{accentOrange}{ÉLEVÉE} & interface\_check \\
  \rowcolor{tableRowA} CISCO-CP-ICMP-002 & Disable ICMP Unreachables & Control & \textcolor{accentGreen}{MOYENNE} & interface\_check \\
  \rowcolor{tableRowB} CISCO-CP-PROXYARP-001 & Disable Proxy ARP & Control & \textcolor{accentGreen}{MOYENNE} & interface\_check \\
  \rowcolor{tableRowA} CISCO-CP-COPP-001 & Control Plane Policing & Control & \textcolor{accentRed}{CRITIQUE} & regex\_present \\
  \rowcolor{tableRowB} CISCO-CP-RACL-001 & Receive ACLs & Control & \textcolor{accentOrange}{ÉLEVÉE} & acl\_check \\
  \rowcolor{tableRowA} CISCO-BGP-GTSM-001 & BGP GTSM TTL Security & Control & \textcolor{accentOrange}{ÉLEVÉE} & regex\_present \\
  \rowcolor{tableRowB} CISCO-BGP-MD5-001 & BGP MD5 Authentication & Control & \textcolor{accentOrange}{ÉLEVÉE} & regex\_present \\
  \rowcolor{tableRowA} CISCO-BGP-MAXPFX-001 & BGP Maximum Prefixes & Control & \textcolor{accentOrange}{ÉLEVÉE} & regex\_present \\
  \rowcolor{tableRowB} CISCO-BGP-PFXLIST-001 & BGP Prefix List Filter & Control & \textcolor{accentOrange}{ÉLEVÉE} & regex\_present \\
  \rowcolor{tableRowA} CISCO-IGP-MD5-001 & IGP MD5 Authentication & Control & \textcolor{accentRed}{CRITIQUE} & regex\_present \\
  \rowcolor{tableRowB} CISCO-IGP-PASSIVE-001 & Passive-Interface & Control & \textcolor{accentGreen}{MOYENNE} & regex\_present \\
  \rowcolor{tableRowA} CISCO-FHRP-AUTH-001 & FHRP Authentication & Control & \textcolor{accentOrange}{ÉLEVÉE} & regex\_present \\
  \rowcolor{tableRowB} CISCO-DP-IPOPT-001 & IP Options Selective Drop & Data & \textcolor{accentGreen}{MOYENNE} & regex\_present \\
  \rowcolor{tableRowA} CISCO-DP-SRCROUTE-001 & Disable IP Source Routing & Data & \textcolor{accentOrange}{ÉLEVÉE} & regex\_present \\
  \rowcolor{tableRowB} CISCO-DP-ICMPRED-001 & Disable ICMP Redirects & Data & \textcolor{accentOrange}{ÉLEVÉE} & interface\_check \\
  \rowcolor{tableRowA} CISCO-DP-DIRCAST-001 & Disable Directed Broadcasts & Data & \textcolor{accentOrange}{ÉLEVÉE} & interface\_check \\
  \rowcolor{tableRowB} CISCO-TACL-ICMP-001 & Transit ICMP Filtering & Data & \textcolor{accentOrange}{ÉLEVÉE} & acl\_check \\
  \rowcolor{tableRowA} CISCO-ANTISPOOFING-URPF-001 & Unicast RPF & Data & \textcolor{accentRed}{CRITIQUE} & interface\_check \\
  \rowcolor{tableRowB} CISCO-ANTISPOOFING-IPSG-001 & IP Source Guard & Data & \textcolor{accentOrange}{ÉLEVÉE} & interface\_check \\
  \rowcolor{tableRowA} CISCO-ANTISPOOFING-DAI-001 & Dynamic ARP Inspection & Data & \textcolor{accentOrange}{ÉLEVÉE} & regex\_present \\
  \rowcolor{tableRowB} CISCO-ANTISPOOFING-PORTSEC-001 & Port Security & Data & \textcolor{accentGreen}{MOYENNE} & interface\_check \\
  \rowcolor{tableRowA} CISCO-TRAFFIC-NETFLOW-001 & NetFlow & Data & \textcolor{accentGreen}{MOYENNE} & regex\_present \\
  \rowcolor{tableRowB} CISCO-VLAN-MAP-001 & VLAN Maps & Data & \textcolor{accentGreen}{MOYENNE} & regex\_present \\
  \rowcolor{tableRowA} CISCO-PVLAN-ISO-001 & Private VLAN Isolated & Data & \textcolor{accentGreen}{MOYENNE} & regex\_present \\
\end{longtable}

% ============================================================
\chapter{Livrables et Planning}
% ============================================================

\section{Livrables Techniques}

\begin{table}[H]
  \centering
  \caption{Livrables du module Cisco}
  \begin{tabularx}{\textwidth}{>{\bfseries}l l X}
    \toprule
    \rowcolor{tableHeader}\color{white}\textbf{Livrable} & \color{white}\textbf{Format} & \color{white}\textbf{Description} \\
    \midrule
    \rowcolor{tableRowA} cisco\_parser.py & Python & Parseur du \texttt{show running-config} \\
    \rowcolor{tableRowB} evaluator.py & Python & Moteur d'évaluation des règles YAML \\
    \rowcolor{tableRowA} management\_plane.yaml & YAML & Règles du Management Plane (~39 règles) \\
    \rowcolor{tableRowB} control\_plane.yaml & YAML & Règles du Control Plane (~12 règles) \\
    \rowcolor{tableRowA} data\_plane.yaml & YAML & Règles du Data Plane (~12 règles) \\
    \rowcolor{tableRowB} reporter.py & Python & Générateur JSON + HTML \\
    \rowcolor{tableRowA} report.json & JSON & Rapport machine-readable \\
    \rowcolor{tableRowB} report.html & HTML & Rapport lisible par l'humain \\
    \rowcolor{tableRowA} Tests unitaires & Python & Couverture $\geq 80\%$ \\
    \bottomrule
  \end{tabularx}
\end{table}

\section{Planning Hebdomadaire Suggéré}

\begin{table}[H]
  \centering
  \caption{Planning de développement}
  \begin{tabularx}{\textwidth}{>{\bfseries\color{ciscoDark}}l X X}
    \toprule
    \rowcolor{tableHeader}\color{white}\textbf{Semaine} & \color{white}\textbf{Objectifs} & \color{white}\textbf{Tâches} \\
    \midrule
    \rowcolor{tableRowA} S1 & Analyse et Architecture & Lecture du guide Cisco, définition du format YAML, prototype du parseur \\
    \rowcolor{tableRowB} S2 & Management Plane & Implémentation des 39 règles Management Plane + tests \\
    \rowcolor{tableRowA} S3 & Control Plane & Implémentation des 12 règles Control Plane (BGP, IGP, CoPP) \\
    \rowcolor{tableRowB} S4 & Data Plane & Implémentation des 12 règles Data Plane (uRPF, DAI, VLAN) \\
    \rowcolor{tableRowA} S5 & Rapport & Générateur HTML/JSON, intégration, démo sur sample \\
    \rowcolor{tableRowB} S6 & Tests \& Finalisation & Tests unitaires, validation sur configs réelles, documentation \\
    \bottomrule
  \end{tabularx}
\end{table}

\section{Critères de Succès}

\begin{itemize}[leftmargin=1.5cm, label=\textcolor{accentGreen}{\faCheckCircle}]
  \item Couverture de \textbf{100\%} des règles définies dans ce cahier des charges.
  \item Taux de faux positifs $< 5\%$ sur un ensemble de configurations de référence.
  \item Génération du rapport en $< 5$ secondes pour un fichier de configuration standard.
  \item Documentation de chaque règle (description, commande IOS, remédiation).
  \item Interface HTML lisible sans connaissances techniques avancées.
\end{itemize}

% ============================================================
\chapter{Annexes}
% ============================================================

\section{Exemple de Fichier de Configuration de Test}

\begin{lstlisting}[caption={Extrait d'un show running-config durcis (exemple)}]
version 16.9
service timestamps debug datetime msec localtime show-timezone
service timestamps log datetime msec localtime show-timezone
service password-encryption
no service tcp-small-servers
no service udp-small-servers
no service finger
no ip source-route
no ip http server
!
hostname ROUTER-CORE-01
!
aaa new-model
aaa authentication login default group tacacs+ local
aaa authorization commands 15 default group tacacs+ local
aaa accounting commands 15 default start-stop group tacacs+
!
enable secret 5 $1$mERr$hx5rVt7rPNoS4wqbXKX7m0
service password-encryption
!
tacacs server TAC-PRIMARY
 address ipv4 10.0.0.100
 key 7 XXXXXXXX
!
ip ssh version 2
ip ssh time-out 60
ip ssh authentication-retries 3
!
ntp authenticate
ntp authentication-key 1 md5 NTPKEY
ntp trusted-key 1
ntp server 10.0.0.200 key 1
!
logging host 10.0.0.150
logging trap informational
logging buffered 64000 informational
logging source-interface Loopback0
no logging console
!
snmp-server view SNMP-VIEW internet included
snmp-server group SNMP-GRP v3 priv read SNMP-VIEW
snmp-server user snmpadmin SNMP-GRP v3 auth sha AUTHPASS priv aes 128 PRIVPASS
!
control-plane
 service-policy input COPP-POLICY
!
banner motd ^
 ACCES REGLEMENTE - ACCES AUTORISE UNIQUEMENT
^
\end{lstlisting}

\section{Exemple de Rapport JSON}

\begin{lstlisting}[language=bash, caption={Structure du rapport JSON}]
{
  "hostname": "ROUTER-CORE-01",
  "assessment_date": "2026-02-19T10:30:00",
  "score": {
    "total_rules": 63,
    "pass": 48,
    "fail": 10,
    "not_found": 3,
    "not_applicable": 2,
    "compliance_pct": 76.2
  },
  "rules": [
    {
      "id": "CISCO-MGT-SSH-001",
      "title": "SSHv2 doit etre active",
      "section": "Management Plane",
      "subsection": "Secure Interactive Management Sessions",
      "status": "PASS",
      "severity": "HIGH",
      "evidence": "ip ssh version 2",
      "remediation": null
    },
    {
      "id": "CISCO-AAA-TYPE7-001",
      "title": "No Type 7 Passwords",
      "section": "Management Plane",
      "subsection": "AAA",
      "status": "FAIL",
      "severity": "CRITICAL",
      "evidence": "tacacs server TAC-PRIMARY key 7 XXXXX",
      "remediation": "Remplacer par type 6 ou key <plaintext> en utilisant 'key 6'"
    }
  ]
}
\end{lstlisting}

\section{Références}

\begin{itemize}[leftmargin=1.5cm]
  \item Cisco Systems, \textit{Harden IOS/IOS-XE Devices}, \url{https://www.cisco.com/c/en/us/support/docs/ip/access-lists/13608-21.html}
  \item Cisco Systems, \textit{NSA Guide to Securing Cisco IOS Devices}
  \item ANSSI, \textit{Recommandations de sécurité pour les architectures réseau}
  \item CIS Cisco IOS XE Benchmark
\end{itemize}

\end{document}
